\section{Results}
\label{sec:results}

\subsection{Interaction Features and Bias}

Table~\ref{tab:statistical_tests} presents the Mann-Whitney U tests comparing interaction features between biased and non-biased articles.
Text length is the most discriminative feature ($+11.5\%$, $p < 10^{-71}$): biased articles tend to be longer.
Among karma features, average comment karma ($+11.1\%$, $p < 10^{-26}$) and karma IQR ($+10.5\%$, $p < 10^{-23}$) show the strongest differences, indicating that biased articles generate more dispersed community reactions.

\begin{table}[t]
\centering
\small
\begin{tabular}{lrrrc}
\toprule
\textbf{Feature} & \textbf{Non-bias} & \textbf{Biased} & \textbf{Diff\%} & \textbf{Sig.} \\
\midrule
text\_length & 343.9 & 383.5 & $+11.5$ & *** \\
avg\_comment\_karma & 3.02 & 3.36 & $+11.1$ & *** \\
karma\_iqr & 2.76 & 3.06 & $+10.5$ & *** \\
pct\_positive\_karma & 0.525 & 0.547 & $+4.1$ & *** \\
median\_comment\_karma & 0.82 & 0.96 & $+16.7$ & *** \\
std\_comment\_karma & 6.62 & 7.24 & $+9.5$ & *** \\
score & 537.1 & 557.8 & $+3.9$ & *** \\
is\_weekend & 0.211 & 0.243 & $+15.5$ & *** \\
avg\_comment\_length & 213.0 & 207.0 & $-2.8$ & *** \\
\bottomrule
\end{tabular}
\caption{Mann-Whitney U tests for top interaction features. Significance: *** $p < 0.001$.}
\label{tab:statistical_tests}
\end{table}

\subsection{Advanced Karma Analysis}

The advanced karma features (Table~\ref{tab:karma_divergence}) reveal that biased articles have significantly higher karma entropy ($+4.6\%$, $p < 10^{-24}$), confirming that their comment sections exhibit more diverse opinion distributions.
The extreme karma fraction (comments beyond $2\sigma$) is also higher ($+2.0\%$, $p < 0.006$), indicating more polarized responses.
Interestingly, the bimodality coefficient and Gini index show no significant differences, suggesting that while biased articles generate more dispersed karma, the shape of the distribution remains structurally similar.

\begin{table}[t]
\centering
\small
\begin{tabular}{lrrrc}
\toprule
\textbf{Feature} & \textbf{Non-bias} & \textbf{Biased} & \textbf{Diff\%} & \textbf{Sig.} \\
\midrule
karma\_entropy & 2.354 & 2.462 & $+4.6$ & *** \\
karma\_q75 & 2.832 & 3.147 & $+11.1$ & *** \\
karma\_iqr & 2.764 & 3.055 & $+10.5$ & *** \\
karma\_extreme\_frac & 0.048 & 0.049 & $+2.0$ & ** \\
karma\_bimodality & 0.783 & 0.789 & $+0.8$ & n.s. \\
karma\_gini & 0.724 & 0.727 & $+0.5$ & n.s. \\
\bottomrule
\end{tabular}
\caption{Advanced karma distribution features. Significance: *** $p < 0.001$, ** $p < 0.01$.}
\label{tab:karma_divergence}
\end{table}

\subsection{Comment Sentiment}

Analysis of 20,000 comments reveals systematic sentiment differences (Table~\ref{tab:sentiment}).
Comments on biased articles are more negative ($59.7\%$ vs.\ $56.2\%$, $p < 10^{-5}$) and contain more anger ($0.271$ vs.\ $0.248$, $p < 4 \times 10^{-5}$).
The polarity score (positive probability minus negative probability) is significantly lower for biased articles ($-0.452$ vs.\ $-0.420$).
Joy, sadness, and fear show no significant differences.

\begin{table}[t]
\centering
\small
\begin{tabular}{lrrrc}
\toprule
\textbf{Metric} & \textbf{Non-bias} & \textbf{Biased} & \textbf{Diff} & \textbf{Sig.} \\
\midrule
\% Negative & 56.2\% & 59.7\% & $+3.5$ & *** \\
\% Positive & 8.0\% & 8.0\% & $0.0$ & n.s. \\
Polarity & $-0.420$ & $-0.452$ & $-0.032$ & *** \\
Anger & 0.248 & 0.271 & $+0.023$ & *** \\
Joy & 0.050 & 0.049 & $-0.001$ & n.s. \\
\bottomrule
\end{tabular}
\caption{Comment sentiment and emotion: biased vs.\ non-biased articles.}
\label{tab:sentiment}
\end{table}

This pattern is remarkably consistent across years.
From 2006 to 2021, comments on biased articles are consistently more negative and carry higher anger scores, with no significant temporal trend in the gap itself.

\subsection{Network Analysis}

\paragraph{Outlet communities.}
Louvain community detection on the outlet projection of the user-outlet bipartite graph identifies two distinct communities (Figure~\ref{fig:communities}, described textually due to space constraints).
Community~1 (26 outlets) includes traditional generalist media (\textit{elmundo.es}, \textit{elpais.com}, \textit{abc.es}, \textit{publico.es}), while Community~2 (24 outlets) comprises digital-native and alternative outlets (\textit{eldiario.es}, \textit{elconfidencial.com}, \textit{xataka.com}).
The most similar outlet pair by shared commenters is \textit{elconfidencial.com}--\textit{eldiario.es} (Jaccard $= 0.40$).

\paragraph{User polarization.}
Of 14,549 active users (${}\geq 10$ comments), 45.4\% fall in the high bias exposure category ($0.6$--$0.8$) and only 8.6\% in the low bias exposure category ($< 0.4$).
Users with mixed exposure read the most diverse media: an average of 33.5 outlets, compared to 12.0 for low-exposure and 12.0 for very-high-exposure users.
This curvilinear relationship between bias exposure and outlet diversity suggests that moderate engagement with biased content correlates with broader information diets.

\subsection{Intra-Article Dynamics}

Table~\ref{tab:dynamics} shows that biased articles trigger faster initial reactions: the time to the 10th comment is $10.7\%$ shorter ($p < 10^{-4}$).
The karma drift (difference between late and early comment karma) is more negative for biased articles ($-8.42$ vs.\ $-7.72$, $p < 10^{-12}$), indicating steeper degradation of discourse quality.
Late comments on biased articles are also shorter ($-4.5\%$, $p < 10^{-9}$) and the total thread lifespan is 20.6\% shorter.

\begin{table}[t]
\centering
\small
\begin{tabular}{lrrrc}
\toprule
\textbf{Feature} & \textbf{Non-bias} & \textbf{Biased} & \textbf{Diff\%} & \textbf{Sig.} \\
\midrule
early\_karma\_mean & 8.19 & 8.98 & $+9.5$ & *** \\
late\_karma\_mean & 0.48 & 0.55 & $+16.6$ & *** \\
karma\_drift & $-7.72$ & $-8.42$ & $-9.1$ & *** \\
karma\_slope & $-0.190$ & $-0.206$ & $-8.4$ & *** \\
first\_10\_hours & 4.26 & 3.80 & $-10.7$ & *** \\
duration\_hours & 127.4 & 101.1 & $-20.6$ & *** \\
late\_comment\_len & 267.4 & 255.4 & $-4.5$ & *** \\
\bottomrule
\end{tabular}
\caption{Intra-article dynamics for articles with $\geq 10$ comments.}
\label{tab:dynamics}
\end{table}

\subsection{Predictive Modeling}

Table~\ref{tab:models} presents classification results.
Gradient Boosting achieves the best performance with an AUC of $0.642$ and macro-F1 of $0.567$, modestly above the majority baseline (accuracy $0.615$).
The most important features for Random Forest are text length ($0.109$), average comment length ($0.079$), and article score ($0.073$).
Logistic Regression coefficients indicate that higher text length ($+0.307$), average comment karma ($+0.253$), and weekend publication ($+0.076$) are positively associated with bias, while higher number of comments ($-0.128$) and comment length ($-0.146$) are negatively associated.

\begin{table}[t]
\centering
\small
\begin{tabular}{lccc}
\toprule
\textbf{Model} & \textbf{Accuracy} & \textbf{F1 Macro} & \textbf{AUC} \\
\midrule
Majority baseline & 0.615 & -- & 0.500 \\
Logistic Regression & 0.618 & 0.487 & 0.607 \\
Random Forest & 0.623 & 0.557 & 0.616 \\
Gradient Boosting & \textbf{0.640} & \textbf{0.567} & \textbf{0.642} \\
\bottomrule
\end{tabular}
\caption{Bias prediction from interaction features (5-fold CV).}
\label{tab:models}
\end{table}

\subsection{Topic and Outlet Analysis}

Bias rates vary substantially across topics.
The highest bias rates appear for interview-related content ($85.2\%$), curiosities ($84.2\%$), and politically charged topics such as ETA ($83.5\%$), gender issues ($81.0\%$), and religion ($80.2\%$).
The lowest rates correspond to technology topics: Microsoft ($16.9\%$), free software ($20.5\%$), and EU policy ($19.4\%$).
Among outlets, \textit{jotdown.es} (88.5\%), \textit{huffingtonpost.es} (87.0\%), and \textit{blogs.publico.es} (85.0\%) have the highest bias rates, while technology-focused outlets show the lowest.
